\subsection{Residual rationality: withdrawal of the claim for 10.35}
\label{cand:10.35}
\problemstatement{10.35}

\paragraph{Correction in version 3.}
The claimed resolution of 10.35 is withdrawn. The Notebook asks about
torsion-free subgroups of \(\GL_n(\overline{\Q})\), with the same
degree \(n\), whereas the argument below concerns \(\GL_2(\Q)\).
Evgeny Khukhro identified the misreading on 19 September. In fact the
displayed torsion-free group already embeds in
\(\GL_2(\Q(i))\leq\GL_2(\overline{\Q})\), so it has the required
residual property. We retain the rational-field result as a separate
observation; it supplies no resolution of the printed question.
Appendix~\ref{app:statement-correction} records the review failure and
the subsequent statement audit.

\paragraph{The rational-field observation.}
Consider
\[
 U(z)=\begin{pmatrix}1&z\\0&1\end{pmatrix},\qquad
 T=\operatorname{diag}(2i,2),\qquad
 G=\langle U(1),T\rangle\leq\GL_2(\C).
\]
Since \(TU(z)T^{-1}=U(iz)\), its unique normal forms are
\(U(z)T^k\), \(z\in\Z[i]\), \(k\in\Z\). Uniqueness follows
from the lower diagonal entry \(2^k\) and then the upper-right
entry. This is \(\Z^2\rtimes\Z\), with the generator acting
by a quarter turn, and is torsion-free because its kernel
\(\Z^2\) and quotient \(\Z\) are torsion-free.

\begin{lemma}
Let \(F\) have characteristic zero and contain no square root
of \(-1\). If \(X,Y\in\GL_2(F)\) satisfy
\[
 [X,YXY^{-1}]=I,\qquad Y^2XY^{-2}=X^{-1},
\]
then \(X^2=I\).
\end{lemma}
\begin{proof}
Assume \(X^2\ne I\). A scalar \(X\) contradicts the second
relation, so \(X\) is nonscalar. Conjugacy to its inverse gives
\(\det(X)^2=1\). If \(\det(X)=-1\), equality of traces gives
\(\tr X=0\), and Cayley--Hamilton gives \(X^2=I\). Thus
\(\det(X)=1\) and \(X^{-1}=(\tr X)I-X\).

The centralizer algebra of a nonscalar \(2\)-by-\(2\) matrix
is \(E=F[I,X]\). Indeed, a vector \(v\) with \(v,Xv\)
independent is cyclic, and a commuting operator is determined
by its value on \(v\). Now \(YXY^{-1}\) is a nonscalar member
of \(E\), so it too generates \(E\) with \(I\).
Conjugation by \(Y\) therefore normalizes \(E\) and fixes
its scalar subspace. On the one-dimensional quotient \(E/FI\)
it acts by a scalar \(q\in F^\times\). Its square sends the
nonzero class of \(X\) to its negative, so \(q^2=-1\),
a contradiction.
\end{proof}

Put \(a=U(1)\). The relations of the lemma hold for \(a,T\)
in \(G\). Consequently every homomorphism
\(\phi:G\to\GL_2(\Q)\), and even every homomorphism into
\(\GL_2(\mathbb R)\), kills the fixed nonidentity element
\(a^2=U(2)\). This disproves residuality over \(\Q\) or \(\R\)
in degree two, but not the property asked for in 10.35.
If the image is torsion-free it kills \(a\) itself, and hence
the whole normal subgroup \(U(\Z[i])\).
The common scalar factor in \(T\) is essential to
torsion-freeness: \(T^4=16I\). The example does embed in
\(\GL_3(\Q)\), so it does not address a version allowing the
target degree to increase. Source:
\artifact{research/10.35-proof.md}.
